\documentclass[11pt]{article}
\usepackage{mathblatt}
% Präambel des Setzers (werkzeuge/setzer-praeambel.tex, wortgleich); dazu \gteil (Muster T2U)
\makeatletter
\setlength{\mbpfbe}{0mm}% keine Punkte (pruefung.md § 5: nur beim Prüfstein)
\newlength{\gEin}\setlength{\gEin}{0pt}
\let\gSchrift\relax
\renewcommand{\pfkreuzzeile}[1]{\par\noindent\hangindent3.2em\hangafter1\hspace*{1.6em}$\square$\ #1\par}
\newcommand{\gkreuz}[1]{$\square$\ #1\hspace{2em}}
\newcommand{\gkreuzzeile}[2]{\par\noindent\hangindent2.2em\hangafter1\mbox{#1\,$\square$}\ #2\par}% Buchstabe, Option: eine je Zeile (neben Skizze)
% links, was man ansieht, rechts, was man tut; Spaltengrenze überall an derselben Stelle (gemeinsam.md „Satz“)
% v0.17: Fußeintrag „Nr: Wert“ – kurz in einer Box (nie getrennt), lang umbrechbar mit der Nummer am Anfang
% (vorher \mbox: lange Lösungen liefen über den Rand, AS40 138 pt); robust, weil \fusshilfe in die .aux schreibt
\NewDocumentCommand{\gfuss}{m m}{\sbox0{#1:~#2}\ifdim\wd0<0.45\textwidth\mbox{#1:~#2}\else\mbox{#1:}\nobreak\ #2\fi}
\newsavebox{\gbox}\newlength{\gLinks}
% v0.10 (Lehrer 13.10.): keine starre Halbierung – die linke Spalte ist so breit wie die Figur (+ 6 mm), die
% Antwort steht direkt daneben; ist die Figur breiter als 0,62 Zeile, steht die Antwort darunter
\newcommand{\gzwei}[2]{\par\smallskip\noindent\sbox{\gbox}{#1}%
  \setlength{\gLinks}{\dimexpr\wd\gbox+6mm\relax}%
  \ifdim\gLinks>0.62\linewidth
    \usebox{\gbox}\par\smallskip\noindent\begin{minipage}{\linewidth}#2\end{minipage}%
  \else
    \begin{minipage}[t]{\gLinks}\vspace{0pt}\usebox{\gbox}\end{minipage}%
    \begin{minipage}[t]{\dimexpr\linewidth-\gLinks\relax}\vspace{0pt}\raggedright #2\end{minipage}%
  \fi\par}
% v0.9 (12.10., pruefung.md § 5 „Marke mit B, Bedingungen rechts“): drittes Argument von pfaufg = Bedingungen
% der Prüfung (v0.10: \pfohnetr = TR durchgestrichen, „FOR“), klein grau im rechten Rand auf der Grundlinie der Nummer, Zeilen mit \\
\newcommand{\gbedrand}[2]{\ifblank{#2}{}{\rlap{\hspace*{\dimexpr#1+2mm\relax}%
  \parbox[t]{15mm}{\raggedright\scriptsize\leavevmode\color{mbgrau}\strut#2}}}}
% Teil eines Bündels (a, b …): Marke links im Rand, Bedingungen rechts im Rand, auf der Zeile des Teils
\newcommand{\gteilrand}[2]{\llap{\mbpfrandmarke{#1}\hspace*{\mbpfnr}}\gbedrand{\linewidth}{#2}}
% Teil eines Bündels: Buchstabe in eigener Spalte, Inhalt eingerückt (\gEin = Einzug, damit die Spaltengrenze von
% \gzwei an derselben Stelle bleibt)
\newenvironment{gbteil}[3]{\par\smallskip\noindent\gteilrand{#2}{#3}\makebox[\mbpfteil][l]{\textbf{#1}}%
  \begin{minipage}[t]{\dimexpr\linewidth-\mbpfteil\relax}\setlength{\gEin}{\mbpfteil}\raggedright\strut\ignorespaces}%
 {\end{minipage}\par}
% Grafik höchstens so breit wie die Spalte (Bündel nebeneinander, Koordinatensystem ohne Karo)
\newsavebox{\gpbox}
\newcommand{\gpassend}[1]{\sbox{\gpbox}{#1}\ifdim\wd\gpbox>\linewidth\typeout{GPASSEND: verkleinert \the\wd\gpbox\space auf \the\linewidth}\resizebox{\linewidth}{!}{\usebox{\gpbox}}\else\usebox{\gpbox}\fi}
% Skizzenfeld (v0.10, Lehrer 13.10.): Karo ab dem linken Rand über die Zeile, „Skizze“ klein grau oben links im Feld
% (wie \pfskizzenfeld der Vorlage); der rechte Rand bleibt den Bedingungen
\newcommand{\gskizze}[1]{\par\smallskip\noindent\begin{tikzpicture}
  \clip (0,0) rectangle ({\linewidth-1mm},{#1*10mm});
  \draw[black!16,line width=0.3pt,step=5mm] (0,0) grid ({\linewidth},{#1*10mm});
  \node[anchor=north west,inner sep=2pt,font=\scriptsize,text=mbgrau,fill=white] at (0,{#1*10mm}) {Skizze};
  \end{tikzpicture}\par}
\RenewDocumentEnvironment{pfaufg}{m m m}{%
  \begin{lrbox}{\mbaufgabebox}\begin{minipage}[t]{\linewidth}%
  \gSchrift\noindent\hspace*{\gEin}\mbpfrandmarke{#1}\gbedrand{\linewidth-\gEin}{#3}%
  \begin{minipage}[t]{\mbpfnr}\strut\textbf{#2}\end{minipage}%
  \begin{minipage}[t]{\dimexpr\linewidth-\gEin-\mbpfnr\relax}\raggedright\strut\ignorespaces}%
 {\end{minipage}%
  \end{minipage}\end{lrbox}%
  \setlength{\mbaufgabehoehe}{\dimexpr\ht\mbaufgabebox+\dp\mbaufgabebox\relax}%
  \ifdefempty{\gkopfnach}{\par\addvspace{5pt}\Needspace*{\mbaufgabehoehe}}{%
    % Stufenkopf, der mit dieser (ersten) Aufgabe zusammenbleibt: Platz für beide verlangen, dann den Kopf setzen
    % (12.10., Quadratische: ein Kopf stand allein unten, die Aufgabe mit Koordinatensystem auf der nächsten Seite)
    % v0.17 (AS40 „Terme deuten“ allein unten): hinter dem Kopf steht das Sprungziel (\special, ein Knoten, nach
    % dem die Seite brechen durfte) – darum \nobreak nach dem Kopf; Vorhalt 56 pt = Kopf mit Linie und Abständen
    \par\Needspace*{\dimexpr\mbaufgabehoehe+56pt\relax}\gkopfnach\par\nobreak\gdef\gkopfnach{}\addvspace{5pt}}%
  \noindent\usebox{\mbaufgabebox}\par\addvspace{5pt}}
% Stufenkopf vormerken: \gkopfvor{<Kopf>} – die nächste pfaufg setzt ihn, sobald beide zusammen Platz haben
\def\gkopfnach{}
\newcommand{\gkopfvor}[1]{\gdef\gkopfnach{#1}}
% Teilaufgabe mit Skizze links in gzwei: Buchstabe in eigener Spalte, darunter Text und Skizze
\newcommand{\gteil}[2]{\makebox[\mbpfteil][l]{\textbf{#1}}\parbox[t]{\dimexpr0.5\textwidth-\gEin-\mbpfnr-5mm-\mbpfteil\relax}{\raggedright #2}}
% Stufenkopf der Vorlage (mathblatt.sty 2026-10-12: Linie, Abstand, 13 pt, \pfjahre) ohne ihren Vorhalt von
% 14 Zeilen: der Setzer fordert vor dem Kopf 0,16 Seitenhöhe (Kopf + erste Aufgabe), mehr nicht
\setlength{\mbpfstufevorhalt}{0pt}
\makeatother
\newcommand{\karo}[1]{\par\smallskip\noindent\begin{tikzpicture}
  \clip (0,0) rectangle ({\linewidth-1mm},{#1*10mm});
  \draw[black!16,line width=0.3pt,step=5mm] (0,0) grid ({\linewidth},{#1*10mm});
  \end{tikzpicture}\par}
\newcommand{\kk}[1]{$\square$\,#1\hspace{0.9em}}
\newcommand{\linie}[1]{\,{\color{black!45}\rule[-1pt]{#1}{0.4pt}}\,}
\newcommand{\pkt}{\textperiodcentered{}}
% v0.14 (14.10.): Stufenkopf mit fünf Punkten (Option = Jahre für die Bestellkarte) oder „zuletzt 19“
\NewDocumentCommand{\gpunkte}{o m}{#2}
\newcommand{\gzuletzt}[1]{zuletzt #1}
% v0.7 (12.10.): Schülerfrage unter dem Stufenkopf; Ankreuzoption als Bild (Buchstabe und Kästchen darunter)
\newcommand{\gfrage}[1]{\par\nobreak\noindent\hspace*{\mbpfnr}{\small\itshape #1}\par\nobreak\addvspace{2pt}}
\newlength{\gbildbreite}\setlength{\gbildbreite}{0.32\linewidth}% Setzer setzt sie je Reihe (vier Bilder: 0,24)
\newcommand{\gbildoption}[3]{\begin{minipage}[t][#3][b]{\gbildbreite}\centering #2\par\smallskip\mbox{#1\,$\square$}\end{minipage}}
% Lineare (12.10.): Bild mit Buchstabe ohne Kästchen (Diagramm zuordnen); Aussage mit „wahr / falsch“ zum Ankreuzen
\newcommand{\gbildname}[3]{\begin{minipage}[t][#3][b]{\gbildbreite}\centering #2\par\smallskip\mbox{\textbf{#1}}\end{minipage}}
\newcommand{\gwahrfalsch}[2]{\par\noindent\hangindent2.2em\hangafter1\makebox[2.2em][l]{(#1)}#2\hfill\mbox{wahr $\square$\quad falsch $\square$}\par}
\begin{document}

\pfheftstil
\pflaufkopf[Abitur GK \pkt{} Prüfung]{Flächeninhalt}{AA150}\pfseitevon
\pfheftkopf{Flächeninhalt}{}
\pfstufenzeile{Fläche gesucht \pkt{} Fläche im Sachzusammenhang gesucht \pkt{} Lösungsweg erklären}

% ===================== Stufe: Fläche gesucht =====================
\gkopfvor{\pfstufekopf{Fläche gesucht}{\gpunkte[26 25 24]{\textcolor{black!80}{$\bullet$}\,\textcolor{black!80}{$\bullet$}\,\textcolor{black!80}{$\bullet$}\,\textcolor{black!25}{$\circ$}\,\textcolor{black!25}{$\circ$}}}\special{pdf:dest (stufe-flache-gesucht) [@thispage /XYZ @xpos @ypos null]}}
% kein Einstieg: keine Schülerfrage in der Gliederung (Schülerfrage, -Fälle, -Form, -Weg, -Skizze; oder Frage) – Einstieg nur von dort, nie aus der Bank (pruefung.md § 5, 12.10.)
% 2025-bebb-gk-A1.4b-m1 (Original 2025-bebb-gk-A1.4b, Schwere II, rang 18; herausgelöste Fassung: keine)
\begin{pfaufg}{GK~’25\\A1.4b}{1.}{\pfohnetr}
Gegeben ist die Funktion $f$ mit $f(x) = 2e^x - 2e^{-1}$ mit der Nullstelle $-1$ (Abbildung). Der Graph von $f$ schließt mit den Koordinatenachsen eine Fläche ein. Berechne ihren Inhalt.
\gzwei{\begin{abiksys}[xmin=-4,xmax=1,ymin=-1,ymax=5,le=0.4]\funktion[0.6]{2*exp(\x)-2*exp(-1)}{}\end{abiksys}}{\karo{2}}
\fusshilfe{\gfuss{1}{Stammfunktion 2$e^{x}$ $-$ 2$e^{ - 1}$ · x; 0,53}}
\end{pfaufg}
% 2024-bebb-gk-A1.1b-m1 (Original 2024-bebb-gk-A1.1b, Schwere II, rang 19; herausgelöste Fassung: keine)
\begin{pfaufg}{GK~’24\\A1.1b}{2.}{\pfohnetr}
Gegeben ist die Funktion $f$ mit $f(x) = x^3 - 9x$. Der Graph von $f$ und die $x$-Achse schließen eine Fläche aus zwei Flächenstücken ein. Berechne ihren Inhalt.
\karo{2}
\fusshilfe{\gfuss{2}{Nullstellen $-$3, 0 und 3; $\tfrac{81}{2}$}}
\end{pfaufg}
% 2026-bb-gk-B2.2d-m1 (Original 2026-bb-gk-B2.2d, Schwere II, rang 20; herausgelöste Fassung: keine)
\begin{pfaufg}{GK~’26\\B2.2d}{3.}{}
Gegeben ist die Funktion $f$ mit $f(x) = 2e^x + 3$. Der Graph von $f$, die Gerade $y = 3$, die $y$-Achse und die Gerade $x = u$ mit $u < 0$ begrenzen eine Fläche. Weise nach, dass ihr Inhalt $2 - 2e^u$ beträgt.
\karo{1}
\fusshilfe{\gfuss{3}{2 $-$ 2$e^{u}$}}
\end{pfaufg}

% ===================== Stufe: Fläche im Sachzusammenhang gesucht =====================
\gkopfvor{\pfstufekopf{Fläche im Sachzusammenhang gesucht}{\gpunkte[26 25 24 22]{\textcolor{black!80}{$\bullet$}\,\textcolor{black!80}{$\bullet$}\,\textcolor{black!80}{$\bullet$}\,\textcolor{black!25}{$\circ$}\,\textcolor{black!80}{$\bullet$}}}\special{pdf:dest (stufe-flache-im-sachzusammenhang-gesucht) [@thispage /XYZ @xpos @ypos null]}}
% kein Einstieg: keine Schülerfrage in der Gliederung (Schülerfrage, -Fälle, -Form, -Weg, -Skizze; oder Frage) – Einstieg nur von dort, nie aus der Bank (pruefung.md § 5, 12.10.)
% 2024-bebb-gk-B2.1j-m1 (Original 2024-bebb-gk-B2.1j, Schwere II, rang 21; herausgelöste Fassung: keine)
\begin{pfaufg}{GK~’24\\B2.1j}{4.}{}
Ein Schrankknauf hat den Querschnitt zwischen den Graphen von $f$ mit $f(x) = (x - 3) \cdot e^{-\tfrac{x}{3} + 2}$ und $-f$ für $3 \le x \le 12$; $F$ mit $F(x) = -3x \cdot e^{-\tfrac{x}{3} + 2}$ ist eine Stammfunktion von $f$. 1 LE entspricht 0,5 cm. Bestimme die Größe der Querschnittsfläche in cm².
\karo{2}
\fusshilfe{\gfuss{4}{9,8 cm²}}
\end{pfaufg}
% 2026-bb-gk-B2.1g-m1 (Original 2026-bb-gk-B2.1g, Schwere II, rang 22; herausgelöste Fassung: keine)
\begin{pfaufg}{GK~’26\\B2.1g}{5.}{}
Eine Modelleisenbahn (Draufsicht, 1 LE = 1 m) ist symmetrisch zu beiden Achsen; oben verläuft sie auf $[-1; 0]$ entlang $G$ mit $g(x) = \tfrac{8}{3}x^3 + 4x^2$, $A\left(-1 \mid \tfrac{4}{3}\right)$, links und rechts schließen Halbkreise um $(\pm 1 \mid 0)$ an. Berechne den Inhalt der eingeschlossenen Fläche in m².
\gzwei{\begin{abiksys}[xmin=-2.6,xmax=2.6,ymin=-1.6,ymax=1.6,le=0.9,xstep=0.5,ystep=0.5]\funktionab[-0.75]{8/3*\x^3+4*\x^2}{G}{-1}{0}\funktionab{-8/3*\x^3+4*\x^2}{}{0}{1}\funktionab{-8/3*\x^3-4*\x^2}{}{-1}{0}\funktionab{8/3*\x^3-4*\x^2}{}{0}{1}\draw[line width=0.8pt] (-1,4/3) arc[start angle=90,end angle=270,radius=4/3];\draw[line width=0.8pt] (1,-4/3) arc[start angle=-90,end angle=90,radius=4/3];\punkt{-1}{4/3}{A}\end{abiksys}}{\karo{2}}
\fusshilfe{\gfuss{5}{$\tfrac{2}{3}$; $\tfrac{16}{9}$ $\pi$; 8,25 m²}}
\end{pfaufg}
% 2022-bebb-gk-B2.2h-m1 (Original 2022-bebb-gk-B2.2h, Schwere III, rang 23; herausgelöste Fassung: keine)
\begin{pfaufg}{GK~’22\\B2.2h}{6.}{}
Ein Betonteil (1 LE = 1 m) hat den Querschnitt $ABDC$: oben der Graph von $h$ mit $h(x) = -\tfrac{1}{16}x^3 + \tfrac{3}{8}x^2$ von $A(0 \mid 0)$ bis $B(4 \mid 2)$, unten die Gerade $g$: $y = \tfrac{1}{2}x - 1$ von $C(0{,}4 \mid -0{,}8)$ bis $D(4 \mid 1)$, seitlich $AC$ und $BD$. Ermittle die Querschnittsfläche (Kontrolle: 3,8 m²).
\karo{2}
\fusshilfe{\gfuss{6}{4; Dreieck unter AC (y = $-$2x von 0 bis 0,4): 0,16; 3,8 m²}}
\end{pfaufg}
% 2025-bebb-gk-B2.1e-m1 (Original 2025-bebb-gk-B2.1e, Schwere III, rang 24; herausgelöste Fassung: keine)
\begin{pfaufg}{GK~’25\\B2.1e}{7.}{}
Ein Windrad: Graph von $f$ mit $f(x) = \tfrac{1}{18}x^3 - \tfrac{3}{2}x$ und Gerade $g$: $y = \tfrac{1}{2}x$ auf $[-6; 6]$, um 90° gedreht (gestrichelt); $A(6 \mid 3)$ liegt auf dem Kreis. Es ist $A_1 = \left(6^2 + 3^2\right) \cdot \pi$ und $A_2 = 4 \cdot \int_{0}^{6} (g(x) - f(x))\,dx$. Beschreibe die Fläche $A_1 - A_2$ in der Abbildung und berechne ihren Inhalt.
\gzwei{\begin{abiksys}[xmin=-7,xmax=7,ymin=-7,ymax=7,le=0.25,xstep=2,ystep=2]\draw[line width=0.6pt] (0,0) circle (6.708);\begin{scope}[rotate=90,dashed]\funktionab{\x^3/18-1.5*\x}{}{-6}{6}\funktionab{\x/2}{}{-6}{6}\end{scope}\funktionab{\x^3/18-1.5*\x}{G_f}{-6}{6}\funktionab{\x/2}{g}{-6}{6}\end{abiksys}}{\karo{2}}
\fusshilfe{\gfuss{7}{Fläche im Kreis außerhalb der Flügel; 69,4}}
\end{pfaufg}

% ===================== Stufe: Lösungsweg erklären =====================
\gkopfvor{\pfstufekopf{Lösungsweg erklären}{\gpunkte[24 23 22]{\textcolor{black!25}{$\circ$}\,\textcolor{black!25}{$\circ$}\,\textcolor{black!80}{$\bullet$}\,\textcolor{black!80}{$\bullet$}\,\textcolor{black!80}{$\bullet$}}}\special{pdf:dest (stufe-losungsweg-erklaren) [@thispage /XYZ @xpos @ypos null]}}
% kein Einstieg: keine Schülerfrage in der Gliederung (Schülerfrage, -Fälle, -Form, -Weg, -Skizze; oder Frage) – Einstieg nur von dort, nie aus der Bank (pruefung.md § 5, 12.10.)
% 2024-bebb-gk-B2.2e-m1 (Original 2024-bebb-gk-B2.2e, Schwere III, rang 25; herausgelöste Fassung: keine)
\begin{pfaufg}{GK~’24\\B2.2e}{8.}{}
Gegeben ist die Funktion $f$ mit $f(x) = 0{,}5x^4 - 3x^2 + 2{,}5$; die Abbildung zeigt einen Ausschnitt ihres Graphen $G_f$. Lösungsweg: (1) $2{,}5x + 2{,}5 = f(x)$ hat für $-1 \le x \le 0$ die Lösungen $-1$ und $0$. (2) $\int_{-1}^{0} (f(x) - (2{,}5x + 2{,}5))\,dx = \tfrac{7}{20}$. (3) $\tfrac{1}{2} \cdot 1 \cdot f(0) = \tfrac{5}{4} = \tfrac{25}{20}$. (4) $7 : 25$. Beschreibe die geometrische Bedeutung jedes Schritts. Zeichne die zugehörigen Flächen in die Abbildung ein.
\par\smallskip\noindent\begin{abiksys}[xmin=-1.5,xmax=1.5,ymin=-1,ymax=3,le=0.55]\funktion[0.9]{0.5*\x^4-3*\x^2+2.5}{G_f}\end{abiksys}\par
\karo{2}
\fusshilfe{\gfuss{8}{(4) Flächenverhältnis 7 : 25}}
\end{pfaufg}
% 2023-bebb-gk-B2.1m-m1 (Original 2023-bebb-gk-B2.1m, Schwere III, rang 26; herausgelöste Fassung: keine)
\begin{pfaufg}{GK~’23\\B2.1m}{9.}{}
Tischplatte wie auf dem Blatt „Integral berechnen“ (Rand im III. Quadranten: Graph von $f$ mit $f(x) = \tfrac{1}{3}(x^2 - 9) \cdot e^{\tfrac{x}{2}}$, 1 LE = 10 cm). Jede Randlinie wird an der Geraden durch die beiden Achsenschnittpunkte ihres Quadranten gespiegelt. Skizziere die neuen Randlinien. Beschreibe einen Lösungsweg für den Zuwachs der Fläche in cm².
\par\smallskip\noindent\begin{abiksys}[xmin=-3.5,xmax=3.5,ymin=-3.5,ymax=3.5,le=0.4]\funktionab[-1.5]{(\x^2-9)*exp(\x/2)/3}{G}{-3}{0}\funktionab{-(\x^2-9)*exp(\x/2)/3}{}{-3}{0}\funktionab{(\x^2-9)*exp(-\x/2)/3}{}{0}{3}\funktionab{-(\x^2-9)*exp(-\x/2)/3}{}{0}{3}\end{abiksys}\par
\karo{2}
\fusshilfe{\gfuss{9}{vier nach außen gewölbte Bögen von Spitze zu Spitze; 687 cm²}}
\end{pfaufg}
% 2022-bebb-gk-B2.1i-m1 (Original 2022-bebb-gk-B2.1i, Schwere III, rang 27; herausgelöste Fassung: keine)
\begin{pfaufg}{GK~’22\\B2.1i}{10.}{}
Gegeben ist die in $\mathbb{R}$ definierte Funktion $f$ mit $f(x) = (x + 3) \cdot e^{-x}$ und $f'(x) = -(x + 2) \cdot e^{-x}$. Im 4. Quadranten schließen der Graph von $f'$, die Koordinatenachsen und die Gerade $x = b$ ($b > 0$) eine Fläche mit dem Inhalt $A_b$ ein. Begründe, dass es ein $b$ gibt mit $A_b = 2$.
\karo{2}
\fusshilfe{\gfuss{10}{3 $-$ (b + 3)$e^{ - b}$; der Wert 2 wird erreicht (b $\approx$ 1,51)}}
\end{pfaufg}

\pfblattende{10}{26 & 25 & 24 & 23 & 22 & 21 & 20 & 19 & 18 & 17}{B2.1g B2.2d & A1.4b B2.1e & A1.1b B2.1j B2.2e & B2.1m & B2.1i B2.2h & A1.1b &  & A1.2b &  & }%
{\pfausschnitt{Integral berechnen}{Flächeninhalt}{Integral am Graphen}}
\end{document}